\documentclass[11pt,letterpaper]{article}
\usepackage[margin=1in]{geometry}
\usepackage{amsmath,amssymb,amsthm}
\usepackage{physics}
\usepackage{hyperref}
\usepackage{booktabs}

\numberwithin{equation}{section}
\renewcommand{\theequation}{S19.\arabic{equation}}

\title{Supplement S19: PMNS--CKM Unification \\
via Bilinear and Vertex Partitions \\
\large Supporting Manuscript \S\S~XV, XVII}
\author{UDT Collaboration}
\date{}

\begin{document}
\maketitle

\begin{abstract}
The six PMNS mixing parameters and the six CKM mixing parameters are
derived from the same Diophantine triple
$(j, \ell, |\kappa_{\max}|) = (1/2, 1, 3)$ through two distinct
partition mechanisms on $S^2$.  PMNS angles arise from bilinear
(two-field) partitions of the quantum numbers, producing large angles.
CKM angles arise from vertex (three-field) partitions, producing small
angles.  The PMNS--CKM hierarchy is a geometric consequence of
partition rank, not a free parameter.
\end{abstract}

\section{The Two Partition Mechanisms}

The angular coupling on $S^2$ admits two classes of quantum number
partition:

\paragraph{Bilinear partition ($\bar\psi\psi$, rank 2).}
Two-field couplings use \emph{squared} quantum numbers.  The partition
fractions are ratios of squared multiplicities:
\begin{equation}
  \sin^2\theta_{12}^{\rm PMNS} = \frac{(2j+1)^2}{(2j+1)^2 + (2\ell+1)^2}
  = \frac{4}{13}\,,
\end{equation}
\begin{equation}
  \sin^2\theta_{23}^{\rm PMNS} = \frac{(2j+1)^2}{2|\kappa_{\max}|+1}
  = \frac{4}{7}\,,
\end{equation}
\begin{equation}
  \sin^2\theta_{13}^{\rm PMNS}
  = \frac{1}{(2\ell+1)^2(2|\kappa_{\max}|-1)} = \frac{1}{45}\,.
\end{equation}
These are $O(0.02$--$0.57)$ --- large angles.

\paragraph{Vertex partition ($\bar\psi A\psi$, rank 3).}
Three-field couplings use \emph{cubed} spin quantum numbers.  The
Cabibbo angle:
\begin{equation}
  \sin\theta_C = \frac{(2\ell+1)^2}{(2j+1)^3(2|\kappa_{\max}|-1)}
  = \frac{9}{40} = 0.225\,,
\end{equation}
and the Wolfenstein parameter:
\begin{equation}
  A = |\phi_0| = \cos(\pi/5) = 0.809\,.
\end{equation}
These are $O(0.04$--$0.23)$ --- small angles.

\section{Why PMNS Is Large and CKM Is Small}

The hierarchy $\theta_{\rm PMNS} \gg \theta_{\rm CKM}$ follows from
partition rank:
\begin{itemize}
\item PMNS: denominators are sums of \emph{two} squared terms
  ($4+9 = 13$, or single terms $7$, $45$).  These are moderate
  numbers, giving large fractions.
\item CKM: denominators involve \emph{products} of cubed spin with
  other multiplicities ($8 \times 5 = 40$).  Products grow faster than
  sums, giving small fractions.
\end{itemize}
This is not a parameter choice --- it is a mathematical property of
partition rank on finite multiplicities.

\section{The CP-Violating Phases}

The PMNS CP~phase is self-referential:
\begin{equation}
  \delta_{\rm CP}^{\rm PMNS} = \pi + \arcsin\frac{4}{13} = 197.9°
  \qquad(0.04\sigma)\,.
\end{equation}
The CKM CP~phase encodes the force hierarchy:
\begin{equation}
  \delta_{\rm CKM} = \arctan\frac{9}{4} = 66.0°
  \qquad(0.1\sigma)\,,
\end{equation}
where $\tan\delta = 9/4 = \alpha_s/\alpha_{\rm EM}$.  CP violation
in the quark sector IS the ratio of the two strongest forces.

\section{Complete Parameter Count}

\begin{center}
\begin{tabular}{lccc}
\toprule
Parameter & Formula & Value & NuFIT/PDG \\
\midrule
$\sin^2\theta_{12}$ (PMNS) & $4/13$ & 0.3077 & $0.304 \pm 0.012$ \\
$\sin^2\theta_{23}$ (PMNS) & $4/7$ & 0.5714 & $0.573 \pm 0.018$ \\
$\sin^2\theta_{13}$ (PMNS) & $1/45$ & 0.02222 & $0.02220 \pm 0.00062$ \\
$\delta_{\rm CP}$ (PMNS) & $\pi + \arcsin(4/13)$ & $197.9°$ & $197° \pm 25°$ \\
$\sin\theta_C$ (CKM) & $9/40$ & 0.2250 & 0.2253 \\
$A$ (Wolfenstein) & $\cos(\pi/5)$ & 0.8090 & 0.811 \\
$\delta_{\rm CKM}$ & $\arctan(9/4)$ & $66.0°$ & $65.5°$ \\
$\bar\rho$ & $1/(2\pi)$ & 0.1592 & 0.159 \\
$\bar\eta$ & $9/(8\pi)$ & 0.3581 & 0.348 \\
$\theta_{\rm QCD}$ & 0 (exact) & 0 & $< 10^{-10}$ \\
\bottomrule
\end{tabular}
\end{center}

Ten mixing parameters from three quantum numbers.  The generation
count $(2\ell+1) = 3$ --- three generations because $\ell = 1$.

\section{The Interconnection}

The same integers $(2,3,5,7)$ that determine the mixing parameters
also determine:
\begin{itemize}
\item The particle masses (Sec.~IX of the manuscript)
\item The fine-structure constant (Sec.~XII, via $36 = 4 \times 9$)
\item The CMB peak count ($7 = 2|\kappa_{\max}|+1$)
\item The nuclear coupling constants (Sec.~XX, via $g_A = 4/\pi$)
\item The He-4 abundance (Sec.~XXV, via $r_b/r_* = 1/\varphi_{\rm gold}$)
\end{itemize}
The mixing parameters are not an independent sector --- they are the
same Diophantine structure viewed through the lens of flavor coupling.

\end{document}
